\section{Related Work}
\label{sec:related_work}

\mypar{Video Diffusion Models.}
Diffusion models~\cite{sohldickstein2015diffusion,ho2020denoising,song2020denoising} have demonstrated amazing capabilities for video generation, conditioned on natural language~\cite{ho2022video,ho2022imagen,guo2023animatediff,girdhar2023emu,bar2024lumiere} or by ``animating'' static images into videos~\cite{singer2022make,blattmann2023stable,xing2024dynamicrafter}.
Beyond content creation~\cite{polyak2024movie}, they can be seen as a path to the ambitious goal of creating world simulators~\cite{videoworldsimulators2024}, showing preliminary success in visual planning for embodied agents~\cite{du2024learning,escontrela2024video,yang2024video}.
Meanwhile, whether the video prior captures sufficient understanding of the physical world is still under debate~\cite{kang2024far}, and explicit integration of physics rules appears to be necessary~\cite{yuan2023physdiff,liu2024physgen,yu2024lucidsim}. 
Our motion prompting technique, applicable to any video diffusion model, not only offers a more flexible and accurate interface to specify motion patterns for video generation, but also serves as a framework to probe %
a trained generative model for their 3D or physics understanding.

\mypar{Motion-conditioned Video Generation.}
A pre-trained text-to-video model can be adapted to follow new motion patterns or additional motion conditioning signals.
Low-rank adaptation (LoRA)~\cite{hu2021lora}, a generic technique for parameter fine-tuning, can be utilized for few-shot motion customization~\cite{ren2024customize,zhao2025motiondirector}.
DreamBooth~\cite{ruiz2022dreambooth}, originally for personalized image generation, can also be applied to video generation~\cite{wu2024motionbooth} with motion control.

Early work proposes video control through sparse motions~\cite{hao2018controllable,ardino2021click}. More recent work explores similar ideas with more powerful models. %
The approaches vary in their design choices but often require certain complicated engineering techniques for stable training and better convergence.
Tora~\cite{zhang2024tora}, MotionCtrl~\cite{wang2024motionctrl}, DragNUWA~\cite{yin2023dragnuwa}, Image Conductor~\cite{li2024image}, and MCDiff~\cite{chen2023mcdiff} adopt two-stage (\eg, finetuning with first dense and then sparse trajectories, or training adapters sequentially), specialized losses~\cite{li2024image,namekata2024sg}, architectures~\cite{xiao2024trajectory,feng2024i2vcontrol}, or multi-stage fine-tuning for multiple modules~\cite{2023videocomposer,chen2023mcdiff,shi2024motion}.
MOFA-Video~\cite{niu2024mofa} requires separate adapters for different motion types, TrackGo~\cite{zhou2024trackgo} uses custom losses and layers, while other works~\cite{yin2023dragnuwa,wang2024motionctrl,niu2024mofa,li2024image,zhang2024tora} engineer data filtering pipelines.
In contrast, we find that a simpler training recipe yields high quality results. Our model is trained in a single stage, with uniformly sampled dense trajectories, and without any specialized engineering efforts. %
Yet it handles a wide range of tasks and motions, generalizing to both sparse and dense trajectories during inference.

Other approaches use entity-centric control signals such as bounding boxes~\cite{wu2024motionbooth,wang2024boximator}, segmentation masks~\cite{wu2024draganything, dai2023finegrained}, human pose~\cite{hu2023animateanyone,xu2023magicanimate}, or camera pose~\cite{he2024cameractrl,watson2024controlling}.
Zero-shot motion adaptation approaches (\eg,~SG-I2V~\cite{namekata2024sg}, Trailblazer~\cite{ma2023trailblazer}, FreeTraj~\cite{qiu2024freetraj}, and Peekaboo~\cite{jain2024peekaboo}) adopt a similar strategy, guiding the video generation based on the motion of entity-centric masks and thus avoiding training or fine-tuning video models.
Our motion prompts offer a more flexible interface to control motion generation at various granularity. Unlike the test-time approaches which explicitly control the diffusion feature maps, our framework naturally balances the strength of controlling signals and that of the encoded video priors.

\mypar{Motion Representations.} 
As our goal is to condition a video generation model on motion of any type, it is crucial to choose a suitable motion representation. 
The most common representation is optical flow~\cite{horn1981determining,lucas1981iterative,brox2010large,dosovitskiy2015flownet,sun2018pwc,teed2020raft}.
While flow can be chained over time, errors can accumulate.
The lack of occlusion handling also makes it unsuitable for our needs, which we find necessary for good camera control (\cref{sec:depth}).
In contrast, long-range feature matching~\cite{jabri2020space,bian2022learning,shrivastava2024selfsupervised,jiang2021cotr} or point trajectories~\cite{harley2022particle,zheng2023point,karaev23cotracker,karaev2024cotracker3,doersch2023tapir,doersch2024bootstap} is a well-suited representation for our application. It can handle occlusions and allows for both sparse and dense tracking over any arbitrary temporal durations. 




